\section{Exact shuffle quotients and a two-product proof at weight five}
\label{gaussian:sec:shuffle}

\subsection{The quotient being computed}

Let \(\mathcal H^1\) be the rational vector space on the empty word
and all binary words ending in \(1\), with shuffle multiplication.
It is bigraded by word length \(w\) and number of ones \(d\).
Let \(\mathcal J\) be its augmentation ideal, spanned by nonempty
words. The indecomposable quotient
\[
Q_{w,d}=(\mathcal J/\mathcal J^2)_{w,d}
\]
sets products of two positive-weight elements equal to zero.
This is a formal combinatorial construction. Shuffle equalities evaluate
to analytic identities at \(z\); the corresponding map on indecomposable
quotients requires also quotienting the evaluated algebra by products.
Evaluation can introduce further relations. In particular, the formal
dimension is not the dimension of a space of numerical periods.

\begin{theorem}[The exact formal shuffle count]\label{gaussian:thm:witt}
For \(w\ge1\) and \(1\le d\le w\),
\[
\dim_{\Q}Q_{w,d}
=\ell_{w,d}
:=\frac1w\sum_{k\mid\gcd(w,d)}
\mu(k)\binom{w/k}{d/k}.
\]
A basis is represented by the binary Lyndon words of length \(w\)
with \(d\) ones. Every such word belongs to \(\mathcal H^1\).
Every other word has a finite, exact polynomial expression in the full
family of Lyndon generators, each monomial having the original total
bidegree. The number of independent product relations in
this component is
\[
\binom{w-1}{d-1}-\ell_{w,d}.
\]
\end{theorem}

\begin{proof}
Order the alphabet by \(0<1\) and words lexicographically.
A Lyndon word is strictly smaller than every nontrivial rotation.
Every word has a unique Chen--Fox--Lyndon factorization into a
nonincreasing concatenation of Lyndon words.
In characteristic zero, shuffle monomials in Lyndon words form a
polynomial basis \cite{gaussian:Radford}. One can see the constructive point
directly: if the factorization is \(u_1^{m_1}\cdots u_r^{m_r}\) with
distinct \(u_1>\cdots>u_r\), then the shuffle product
\(u_1^{\shuffle m_1}\shuffle\cdots\shuffle u_r^{\shuffle m_r}\)
has this concatenation as its lexicographically largest word, with
coefficient \(m_1!\cdots m_r!\). All other words are smaller.
Induction therefore gives a triangular polynomial expression, and the
same triangularity proves linear independence of the monomials.

Every Lyndon word of length greater than one that contains both letters
starts with \(0\) and ends with \(1\). The single-letter words are
\(0\) and \(1\). The factorization of a word ending in \(1\) cannot
contain the factor \(0\), because that smallest factor would be last,
forcing the concatenation to end in \(0\).
Thus \(\mathcal H^1\) is exactly the polynomial algebra on all Lyndon
words other than \(0\). Passing to its augmentation quotient leaves
one class for each generator.

It remains to count. A primitive binary word of length \(w\) with
\(d\) ones has \(w\) distinct rotations and exactly one Lyndon rotation.
Every binary word is a power of a unique primitive word.
M\"obius inversion on this power decomposition gives
\(\sum_{k\mid\gcd(w,d)}\mu(k)\binom{w/k}{d/k}\)
primitive words, proving the formula. Finally, there are
\(\binom{w-1}{d-1}\) words ending in \(1\) in the component, which
gives the claimed product rank.
\end{proof}

\begin{table}[htbp]
\centering
\begin{tabular}{c|rrrrr}
\toprule
Weight \(w\)&\(d=1\)&\(d=2\)&\(d=3\)&\(d=4\)&\(d=5\)\\
\midrule
3&1&1&0&--&--\\
4&1&1&1&0&--\\
5&1&2&2&1&0\\
6&1&2&3&2&1\\
7&1&3&5&5&3\\
8&1&3&7&8&7\\
9&1&4&9&14&14\\
\bottomrule
\end{tabular}
\caption{The exact formal indecomposable counts \(\ell_{w,d}\).
They count Lyndon generators, not independent numerical constants.}
\label{gaussian:tab:witt}
\end{table}

For \(w=6,d=3\), the ten words reduce modulo products to three Lyndon
words:
\[
000111,\qquad001011,\qquad001101,
\]
corresponding respectively to \(F_{4,1,1}\), \(F_{3,2,1}\),
and \(F_{3,1,2}\). This explains the manuscript's finite rank-seven
observation and gives its uniform extension.
The companion code records the full polynomial normal forms of all ten
triples. For a compact view of the quotient, put
\(X=[F_{4,1,1}],Y=[F_{3,2,1}],Z=[F_{3,1,2}]\), where brackets mean
classes modulo products. Exact triangular elimination gives
\[
\begin{array}{c|rrr@{\qquad}c|rrr}
 &X&Y&Z&&X&Y&Z\\ \hline
{}[F_{4,1,1}]&1&0&0 &[F_{1,4,1}]&6&3&1\\
{}[F_{3,2,1}]&0&1&0 &[F_{1,3,2}]&-12&-6&-2\\
{}[F_{3,1,2}]&0&0&1 &[F_{1,2,3}]&12&6&2\\
{}[F_{2,3,1}]&-9&-4&-1 &[F_{1,1,4}]&-6&-3&-1\\
{}[F_{2,2,2}]&12&4&0 &[F_{2,1,3}]&-3&-1&0
\end{array}
\]
For example, one full identity before taking the quotient is
\begin{equation}
F_{2,2,2}(z)=12F_{4,1,1}(z)+4F_{3,2,1}(z)
-2\Li_2(z)F_{3,1}(z)+\frac{\Li_2(z)^3}{6}.
\label{gaussian:eq:full-222}
\end{equation}
Re-expanding its right side into words gives the single word \(010101\)
with coefficient one. This supplies a finite rational certificate,
including the product terms that a quotient table suppresses.

\subsection{An elementary all-weight double elimination}

At depth two the count simplifies to
\[
\ell_{w,2}=\left\lfloor\frac{w-1}{2}\right\rfloor.
\]
There is also a direct triangular proof that avoids general shuffle
algebra. For \(p+q=w\), shuffling the two single-polylogarithm words gives
\begin{align}
\Li_p(z)\Li_q(z)
={}&\sum_{j=1}^{p}\binom{w-j-1}{q-1}F_{w-j,j}(z)
+\sum_{j=1}^{q}\binom{w-j-1}{p-1}F_{w-j,j}(z).
\label{gaussian:eq:general-double-shuffle}
\end{align}
The binomial coefficient counts the interleavings of the zero blocks
before the second nonzero letter.

\begin{proposition}[Independent double product rows]\label{gaussian:prop:double-rank}
For fixed \(w\), the rows of
\eqref{gaussian:eq:general-double-shuffle} with
\(1\le p\le\lfloor w/2\rfloor\) have rank
\(\lfloor w/2\rfloor\) on the \(w-1\) double coordinates.
They express the entire family through products of singles and the
\(\lfloor(w-1)/2\rfloor\) coordinates
\[
\{F_{w-j,j}:1\le j\le\lfloor(w-1)/2\rfloor\}.
\]
This set is empty when \(w=2\).
\end{proposition}

\begin{proof}
In the row indexed by \(p\), the largest inner index appearing is
\(q=w-p\). Its coefficient is \(1\) when \(p<q\), and \(2\)
when \(p=q\). No larger inner index occurs.
Order the rows with \(p\) decreasing and use the columns with
\(j=w-p\); the corresponding square block is triangular with nonzero
diagonal. All products of two depth-one words have already been included,
since swapping \(p,q\) repeats a row. This proves the rank and the
elimination statement.
\end{proof}

At \(z=i\), taking imaginary parts of these equations gives an
unconditional spanning bound for the \(g_{a,b}\), with explicit
single-value terms. It can be sharpened in the parity sector, as the
next section shows.

\subsection{A normalized Pascal inverse for the shuffle rows}
The triangular elimination can be made explicit in outer-index coordinates.
Let $m=\lfloor w/2\rfloor$ and let $A_{p,a}=\binom{a-1}{p-1}+
\binom{a-1}{w-p-1}$, for $1\le p\le m$, $1\le a<w$.
Divide row $p$ by $c_p=1+\mathbf1_{2p=w}$. Its first $m$ columns are
$P_{p,a}=\binom{a-1}{p-1}$, with inverse
\[
 (P^{-1})_{i,j}=(-1)^{j-i}\binom{j-1}{i-1}\quad(j\ge i),
 \qquad (P^{-1})_{i,j}=0\quad(j<i).
\]
Indeed the matrix product reduces to
$\binom{j-1}{i-1}\sum_{k=i}^j(-1)^{j-k}\binom{j-i}{k-i}=\delta_{ij}$.
Writing $R_p=\Imag(\Li_p(\iu)\Li_{w-p}(\iu))$, the constructive solution is
\begin{equation}\label{signed:eq:shuffle-solve}
 g_{a,w-a}=\sum_{p=a}^m(-1)^{p-a}\binom{p-1}{a-1}
 \left[\frac{R_p}{c_p}-\sum_{j=m+1}^{w-1}\frac{A_{p,j}}{c_p}g_{j,w-j}\right],
 \qquad 1\le a\le m.
\end{equation}
This explicitly leaves the same $\lfloor(w-1)/2\rfloor$ formal coordinates
as Proposition~\ref{gaussian:prop:double-rank}; evaluation can impose further
relations.

\subsection{The weight-five relation is two shuffles}

The words for \(\Li_1\Li_4\) and \(\Li_2\Li_3\) give
\begin{align}
\Li_1(z)\Li_4(z)
&=2F_{4,1}(z)+F_{3,2}(z)+F_{2,3}(z)+F_{1,4}(z),
\label{gaussian:eq:shuffle14}\\
\Li_2(z)\Li_3(z)
&=6F_{4,1}(z)+3F_{3,2}(z)+F_{2,3}(z).
\label{gaussian:eq:shuffle23}
\end{align}
The exact single values are
\[
\Li_1(i)=-\frac L2+\frac{\pi i}{4},\quad
\Li_2(i)=-\frac{\pi^2}{48}+iG,\quad
\Li_3(i)=-\frac{3\zeta(3)}{32}+\frac{\pi^3i}{32},
\]
\[
\Li_4(i)=-\frac{7\pi^4}{11520}+i\beta(4).
\]
Consequently
\begin{align}
2g_{41}+g_{32}+g_{23}+g_{14}
&=-\frac12\beta(4)L-\frac{7\pi^5}{46080},
\label{gaussian:eq:wt5-first}\\
6g_{41}+3g_{32}+g_{23}
&=-\frac{\pi^5}{1536}-\frac{3G\zeta(3)}{32}.
\label{gaussian:eq:wt5-second}
\end{align}
Here \(g_{41}\) means \(g_{4,1}\), and similarly for other subscripts.

\begin{corollary}[Proof and strengthening of the recorded relation]
\label{gaussian:cor:wt5}
The experimentally discovered identity
\begin{equation}\label{gauss:eq:wt5-sporadic}
576g_{41}+288g_{32}+736g_{23}+960g_{14}
=21G\zeta(3)-480\beta(4)L
\end{equation}
holds. Moreover the four \(g\)-values are spanned by \(g_{41},g_{32}\)
and the single-value products \(\pi^5,G\zeta(3),\beta(4)L\).
Explicitly,
\begin{align}
g_{23}
&=-6g_{41}-3g_{32}-\frac{\pi^5}{1536}
  -\frac{3G\zeta(3)}{32},\label{gaussian:eq:g23-elim}\\
g_{14}
&=4g_{41}+2g_{32}+\frac{23\pi^5}{46080}
  +\frac{3G\zeta(3)}{32}-\frac12\beta(4)L.
\label{gaussian:eq:g14-elim}
\end{align}
\end{corollary}

\begin{proof}
Multiply \eqref{gaussian:eq:wt5-first} by \(960\), multiply
\eqref{gaussian:eq:wt5-second} by \(-224\), and add.
The \(\pi^5\) coefficients cancel exactly. Solving the two
independent equations instead gives \eqref{gaussian:eq:g23-elim} and
\eqref{gaussian:eq:g14-elim}.
\end{proof}

This proof explains why a search can return an apparently complicated
relation while missing the simpler structure: eliminating a chosen
single-value coordinate can enlarge coefficients.
The mathematical content is visible before the elimination.

\subsection{An infinite odd-weight shuffle family}
\begin{theorem}[An infinite odd-weight family with the pure $\pi$ term removed]
\label{cert:thm:oddfamily}
For $m\ge2$, set $W=2m+1$ and
\begin{align}
 A_W&=2g_{W-1,1}+\sum_{a=1}^{W-2}g_{a,W-a},\notag\\
 B_W&=2(W-2)g_{W-1,1}+(W-2)g_{W-2,2}
       +\sum_{a=2}^{W-3}(a-1)g_{a,W-a},\label{cert:eq:AB}\\
 \kappa_m&=\frac{12(2^{2m-1}-1)B_{2m}}
                   {(2m)(2m-1)E_{2m-2}}.\label{cert:eq:lambda}
\end{align}
Then $\kappa_m$ is a positive rational number and
\begin{equation}
 A_W-\kappa_m B_W
 =-\frac12\beta(2m)L+
 \kappa_m 2^{-(2m-1)}(1-2^{2-2m})G\zeta(2m-1).
\label{cert:eq:odd-family}
\end{equation}
\end{theorem}
\begin{proof}
Equation \eqref{gaussian:eq:general-double-shuffle} identifies $A_W$ and $B_W$ with
$\Imag(\Li_1(\iu)\Li_{2m}(\iu))$ and
$\Imag(\Li_2(\iu)\Li_{2m-1}(\iu))$ respectively.
Writing $c_k=2^{-k}(1-2^{1-k})$, these products are
\begin{align*}
 A_W&=-\tfrac12 L\beta(2m)-\tfrac\pi4 c_{2m}\zeta(2m),\\
 B_W&=-c_{2m-1}G\zeta(2m-1)-\tfrac{\pi^2}{48}\beta(2m-1).
\end{align*}
Their pure-$\pi^W$ terms have ratio
\[
 \frac{12c_{2m}\zeta(2m)}{\pi\beta(2m-1)}.
\]
Euler's even-zeta formula and the classical odd-beta evaluation turn this ratio into
\eqref{cert:eq:lambda}. The signs of $B_{2m}$ and $E_{2m-2}$ agree, so the
ratio is positive. Subtraction proves \eqref{cert:eq:odd-family}.
\end{proof}

The first three coefficients are
\[
 \kappa_2=\frac7{30},\qquad
 \kappa_3=\frac{31}{525},\qquad
 \kappa_4=\frac{127}{8540}.
\]
For example, at weight seven the identity can be written without
fractions as
\begin{equation}
 268800A_7-15872B_7=465G\zeta(5)-134400\beta(6)L.
\label{cert:eq:odd7}
\end{equation}
This theorem does not reduce every odd-weight $g_{a,b}$ individually.
It supplies one explicit family of linear combinations in the parity
component not individually reduced by Theorem~\ref{gaussian:thm:binomial-parity}.


At weight seven a second combination complements the pure-$\pi$-free
family above:
\begin{align}\label{signed:eq:w7-second}
 430g_{6,1}+215g_{5,2}+79g_{4,3}+11g_{3,4}-7g_{2,5}
 =\frac{105}{512}G\zeta(5)-\frac{75}{32}\beta(4)\zeta(3).
\end{align}
Its exact certificate is $-7$ times the $(2,5)$ shuffle row plus $25$ times
the $(3,4)$ row, after evaluating their single values. The exact row calculation fixes the homogeneous weight-seven products.

\section{The mixed Gaussian weight-five identity and its exact proof}
\label{s4proof:sec:s4-proof}

Put $G=\beta(2)$ and
\[
 g_{a,b}=\operatorname{Im}\operatorname{Li}_{a,b}(i,1),
 \qquad
 S_p=\sum_{n=0}^{\infty}\frac{(-1)^nH_n}{(2n+1)^p}.
\]
The manuscript's equation \texttt{gauss:eq:S4-closed} was explicitly
classified as a numerically supported identity. We prove precisely that
identity. The contribution is a reduction of its mixed Gaussian term to
three specified single-color doubles and elementary products. It is not a
claim that the remaining doubles are independent or irreducible, nor a
claim that $S_4$ previously lacked every representation in special functions.
Indeed the 2021 discussion by Hintze and pisco already records the
representation by two colored double values and relates it to
hypergeometric evaluations \cite{s4proof:hintze2021s4,s4proof:mhzhao2020}.

\begin{theorem}[Certified mixed Gaussian reduction]\label{s4proof:thm:s4}\label{gaussian:conj:S4}\label{gaussian:thm:S4}
With the series convention
$\operatorname{Li}_{a,b}(x,y)=\sum_{n>m\ge1}x^ny^m/(n^am^b)$,
\begin{align}
 S_4={}&\frac{4g_{4,1}-3g_{3,2}-9g_{2,3}}7
       +\frac{\pi^5}{224}
       -\frac{27}{224}G\zeta(3)-2\beta(4)\log2.
 \label{gauss:eq:S4-closed}
\end{align}
Equivalently,
\begin{align}
 \operatorname{Im}\operatorname{Li}_{4,1}(i,-1)
 ={}&-\frac{3g_{4,1}+3g_{3,2}+9g_{2,3}}7
       +\frac{\pi^5}{224}
       -\frac{27}{224}G\zeta(3)-2\beta(4)\log2.
 \label{s4proof:eq:s4-mixed-proved}
\end{align}
The proof consists of convergent changes of variables, standard
shuffle/stuffle identities with the explicitly described single-divergence
regularization, and a finite rational linear combination supplied as an
independently replayable certificate.
\end{theorem}

\subsection{The arithmetic translation}
For $n=2k+1$,
\[
 H_k=\sum_{m=1}^{n-1}\frac{1+(-1)^m}{m}.
\]
Taking imaginary parts removes even $n$, while
$\operatorname{Im}(i^{2k+1})=(-1)^k$. Consequently, for $p\ge1$,
\begin{equation}\label{gaussian:eq:S-mixed}
 S_p=\operatorname{Im}\left(
 \operatorname{Li}_{p,1}(i,1)+\operatorname{Li}_{p,1}(i,-1)\right).
\end{equation}
For $p>1$ the interchange is justified by
$\sum_{n\ge2}H_{n-1}/n^p<\infty$. At $p=1$, use the common radial
Abel limit and the ordinary convergence proved by summation by parts.
Thus \eqref{gauss:eq:S4-closed} and \eqref{s4proof:eq:s4-mixed-proved} are equivalent.

\subsection{A convergent octahedral change of variables}
Write $\omega_a=dt/(t-a)$ and
\[
 \mathcal I(a_1\cdots a_w)
 =\int_{0<t_w<\cdots<t_1<1}
     \omega_{a_1}(t_1)\cdots\omega_{a_w}(t_w),
\]
extended multilinearly in the letters. A convergent word has first letter
different from $1$ and last letter different from $0$.
For colors $c_j\in\mathbb Z/4\mathbb Z$, set
\[
 Z((s_1,c_1),\ldots,(s_d,c_d))
 =\operatorname{Li}_{s_1,\ldots,s_d}(i^{c_1},\ldots,i^{c_d}).
\]
Its integral word is
\begin{equation}\label{s4proof:eq:s4-word-convention}
 Z(\boldsymbol s;\boldsymbol c)
 =(-1)^d\mathcal I\left(
 0^{s_1-1}i^{-c_1}\,
 0^{s_2-1}i^{-(c_1+c_2)}\cdots
 0^{s_d-1}i^{-(c_1+\cdots+c_d)}\right).
\end{equation}
This fixes both the color convention and every depth-dependent sign.

Consider the involution
\[
 \phi(t)=\frac{1-t}{1+t}.
\]
It preserves the puncture set $\{0,\infty,1,-1,i,-i\}$ and reverses the
real interval. Direct differentiation gives
\[
 \phi^*\omega_a=\omega_{\phi(a)}-\omega_{-1}\quad(a\ne-1),
 \qquad \phi^*\omega_{-1}=-\omega_{-1},
\]
where $\phi(0)=1$, $\phi(1)=0$, $\phi(i)=-i$, and
$\phi(-i)=i$. Denote this linear map on letters by $\tau$.
Changing variables and reversing the order of the simplex proves
\begin{equation}\label{s4proof:eq:s4-convergent-duality}
 \mathcal I(a_1\cdots a_w)
 =(-1)^w\mathcal I\bigl(\tau(a_w)\cdots\tau(a_1)\bigr).
\end{equation}
No tangential-basepoint correction occurs here: an original last letter
other than $0$ transforms into first letters other than $1$, and an original
first letter other than $1$ transforms into last letters other than $0$.
Every term of the expansion therefore remains convergent.
The broader use of octahedral symmetries to supplement standard level-four
relations goes back to Zhao \cite{s4proof:jqzhao2008,gaussian:Zhao}; the elementary
involution above is sufficient for the present certificate.

Only one instance of \eqref{s4proof:eq:s4-convergent-duality} is needed. The word of
$\operatorname{Li}_{1,4}(i,1)$ is $(-i)000(-i)$, and hence
\begin{equation}\label{s4proof:eq:s4-duality-seed}
 D:=\operatorname{Im}\left\{
 \mathcal I(e_{-i}e_0^3e_{-i})+
 \mathcal I\bigl((e_i-e_{-1})(e_1-e_{-1})^3(e_i-e_{-1})\bigr)\right\}=0.
\end{equation}
Here $e_a$ means the letter at $a$, multiplication is concatenation, and
powers are powers in the noncommutative word algebra. The transformed product expands into
exactly $32$ integral words, with coefficients $\pm1$. The first integral is
the single-color double; the second group is converted to colored
polylogarithms using \eqref{s4proof:eq:s4-word-convention}.

\subsection{The standard relation rows}
The certificate uses two schemas, interpreted over $\mathbb Q$ before
imaginary parts are taken.
First, for two admissible multiple-polylogarithm index/color lists $u,v$,
let $\mathsf H(u,v)$ be the integral shuffle product and
$\mathsf S(u,v)$ the series stuffle product. Then
\begin{equation}\label{s4proof:eq:s4-ds-row}
 R_{u,v}:=\operatorname{Im}\bigl(\mathsf H(u,v)-\mathsf S(u,v)\bigr)=0.
\end{equation}
The same row is used when $u=((1,0))$ and $v$ is admissible, with the standard
single-divergence regularization. This point must not be suppressed:
$\operatorname{Li}_1(1)$ itself is divergent. Its shuffle and stuffle
regularizations are both denoted by a formal variable $T$. A single leading
$(1,0)$ gives a polynomial of degree at most one in $T$, and the
regularization comparison is the identity in degrees zero and one.
Explicitly, its comparison operator satisfies
\[
 \rho(e^{Tu})=\exp\!\left(\sum_{n\ge2}\frac{(-1)^n\zeta(n)}n u^n\right)e^{Tu},
\]
so $\rho(1)=1$ and $\rho(T)=T$ \cite[Theorem~2.9]{gaussian:Zhao}.
Subtracting the two product expansions cancels their common divergent word
and leaves the convergent identity \eqref{s4proof:eq:s4-ds-row}. This is the
single-divergence case of regularized double shuffle
\cite{s4proof:racinet2002,gaussian:Zhao}. The verifier rejects all other divergent input
patterns and checks that every surviving word is admissible.

Second, for a list $z=((s_1,c_1),\ldots,(s_d,c_d))$ with
$c_j\in\{0,1\}$ and with all terms below convergent, ordinary distribution
is
\begin{equation}\label{s4proof:eq:s4-distribution-row}
 Z(\boldsymbol s;2\boldsymbol c)
 -2^{w-d}\sum_{\boldsymbol\epsilon\in\{0,2\}^d}
 Z(\boldsymbol s;\boldsymbol c+\boldsymbol\epsilon)=0,
 \qquad w=s_1+\cdots+s_d.
\end{equation}
Indeed summing signs forces every nested summation index to be even, so the
factor is exactly $2^{w-d}$. The certificate multiplies some such relations
by an admissible value of complementary weight and expands that product by
integral shuffle. These are \emph{lifted convergent distribution} rows; no
regularized distribution formula is assumed. At weight one this includes
the elementary seed $\operatorname{Li}_1(-1)=
\operatorname{Li}_1(i)+\operatorname{Li}_1(-i)$.

\subsection{The exact certificate identity}
In the formal real vector space whose generators are imaginary parts of
convergent colored values of weight five, define
\begin{align}
 F={}&1120\operatorname{Im}\operatorname{Li}_{4,1}(i,-1)
       +480g_{4,1}+480g_{3,2}+1440g_{2,3}
       -1536\operatorname{Im}\operatorname{Li}_5(i)\nonumber\\
 &\quad+135\operatorname{Im}\bigl(\operatorname{Li}_2(i)
                                      \operatorname{Li}_3(1)\bigr)
       -2240\operatorname{Im}\bigl(\operatorname{Li}_4(i)
                                      \operatorname{Li}_1(-1)\bigr).
 \label{s4proof:eq:s4-certificate-target}
\end{align}
The two products in this expression are expanded by integral shuffle.
Complex conjugation is imposed only by
$\operatorname{Im}Z(\boldsymbol s;-\boldsymbol c)
=-\operatorname{Im}Z(\boldsymbol s;\boldsymbol c)$.

\begin{lemma}[Finite rational certificate]\label{s4proof:lem:s4-certificate}
The accompanying file \texttt{S4\_certificate.json} gives $911$ rational
numbers $q_r$ and $911$ standard-relation descriptors $R_r$ such that the
following is an exact equality of formal vectors:
\begin{equation}\label{s4proof:eq:s4-certificate-equality}
 F=-1120D+\sum_{r=1}^{911}q_rR_r.
\end{equation}
The rows consist of $713$ convergent double-shuffle rows, $181$ rows with
one regularized weight-one factor, and $17$ lifted convergent distribution
rows.
\end{lemma}
\begin{proof}
The supplied program \texttt{verify\_s4.py} rebuilds every row from the
indices and colors in its descriptor. Integral shuffles are generated by
choosing interleaving positions, series stuffles by the three defining
recursive cases, and distribution by sign enumeration. The program also
reconstructs the single duality seed and the target independently of their
stored expansions. All coefficients are integers or instances of Python's
exact \texttt{Fraction} type. Summing the left side minus the right side of
\eqref{s4proof:eq:s4-certificate-equality} yields the empty coefficient dictionary,
so every formal-word coefficient is exactly zero. The verifier neither
reruns a rank calculation nor uses numerical polylogarithms, integer-relation
searches, cached matrices, or modular arithmetic. Thus the certificate is
a finite rational proof of the stated formal equality.
\end{proof}

\begin{proof}[Proof of Theorem~\ref{s4proof:thm:s4}]
Every $R_r$ vanishes by \eqref{s4proof:eq:s4-ds-row} or
\eqref{s4proof:eq:s4-distribution-row}, and $D=0$ by
\eqref{s4proof:eq:s4-duality-seed}. The certificate therefore gives $F=0$.
Now
\[
 \operatorname{Im}\operatorname{Li}_5(i)=\beta(5)=\frac{5\pi^5}{1536},
 \qquad \operatorname{Li}_3(1)=\zeta(3),
 \qquad \operatorname{Li}_1(-1)=-\log2.
\]
Substituting these evaluations into \eqref{s4proof:eq:s4-certificate-target} gives
\[
 1120\operatorname{Im}\operatorname{Li}_{4,1}(i,-1)
 +480g_{4,1}+480g_{3,2}+1440g_{2,3}
 -5\pi^5+135G\zeta(3)+2240\beta(4)\log2=0.
\]
Division by $1120$ proves \eqref{s4proof:eq:s4-mixed-proved}, and
\eqref{gaussian:eq:S-mixed} proves \eqref{gauss:eq:S4-closed}.
\end{proof}

\paragraph{Independent numerical audit.}
An independent real-quadrature evaluation uses
\[
 S_4=\frac16\int_0^1\frac{\log^3x\,\log(1+x^2)}{1+x^2}\,dx,
 \qquad
 g_{a,b}=\frac1{(a-1)!}\int_0^1(-\log x)^{a-1}
 \operatorname{Im}\frac{i\operatorname{Li}_b(ix)}{1-ix}\,dx.
\]
At $70$ decimal working digits it gives
\[
 S_4=-0.0104934562973350434533567510062017519635106188773197615131294,
\]
with discrepancy about $4.6\times10^{-71}$ between the two sides.
This numerical agreement is an audit of conventions, not a step in the
proof; the proof is the exact certificate above.

\paragraph{Scope and further work.}
The existing two relations among $g_{4,1},g_{3,2},g_{2,3},g_{1,4}$ may be
used to shorten \eqref{gauss:eq:S4-closed}; they are prior results and are not
counted as new here. The present proof establishes an upper bound on a
specific constant's expression complexity. It says nothing about a lower
bound, linear independence, or the impossibility of a classical-polylogarithm
reduction. A useful next problem is to compress the $911$ standard rows
into a short conceptual derivation, and to apply the same convergent
involution to the analogous odd-weight mixed sums $S_6,S_8,\ldots$.

\begin{corollary}[The two-coordinate form]\label{s4proof:cor:short}
\begin{equation}\label{gaussian:eq:S4-short}
 S_4=\frac{58g_{4,1}+24g_{3,2}}7+\frac{19\pi^5}{3584}-2\beta(4)\log2.
\end{equation}
\end{corollary}
\begin{proof}
Use the two weight-five shuffle eliminations
\eqref{gaussian:eq:g23-elim}--\eqref{gaussian:eq:g14-elim} in
Theorem~\ref{gaussian:thm:S4}. No independence assumption enters.
\end{proof}


\section{An explicit binomial parity reduction for every double weight}
\label{gaussian:sec:double-parity}

Throughout this section the convention is
\[
 F_{a,b}(z)=\Li_{a,b}(z,1)
   =\sum_{n>m>0}\frac{z^n}{n^a m^b},\qquad
 U_n(z)=\operatorname{Re}\Li_n(z),\quad
 V_n(z)=\operatorname{Im}\Li_n(z).
\]
We take $|z|=1$, $z\ne1$, and use radial boundary values. These are also
the ordinary convergent series values: summation by parts applies to
$H_{n-1}^{(b)}/n^a$, whose total variation is finite and which tends to zero.
The result below is an elementary explicit specialization of the established
multiple-polylogarithm parity theorem. It should not be presented as a new
general parity theorem: Panzer proved the general theorem in 2017, including
explicit formulas at depths two and three; Umezawa supplied a general-depth
explicit formula in 2025 \cite{Panzer2017,gaussian:Umezawa}. Its contribution here
is a short independently checkable formula in the manuscript's convention,
including the boundary cases with an index equal to one.

For $w\ge2$ define the parity component
\[
 E_{a,b}(z)=
 \begin{cases}
 \operatorname{Im}F_{a,b}(z),&w=a+b\text{ even},\\
 \operatorname{Re}F_{a,b}(z),&w=a+b\text{ odd},
 \end{cases}
\quad
 W_n(z)=
 \begin{cases}V_n(z),&w\text{ even},\\U_n(z),&w\text{ odd}.
 \end{cases}
\]
The choice of $W_n$ depends on the fixed total weight $w$, not on $n$.

\begin{theorem}[Explicit double parity formula]\label{gaussian:thm:binomial-parity}
Let $w\ge2$, $1\le p\le w-1$, and $q=w-p$. Put
\begin{align*}
 A_p(z)&=\sum_{j=1}^{p}\binom{w-j-1}{q-1}E_{w-j,j}(z),\\
 P_p(z)&=
 \begin{cases}
 U_q(z)V_p(z),&w\text{ even},\ p\text{ odd},\\
 U_p(z)V_q(z),&w\text{ even},\ p\text{ even},\\
 -V_p(z)V_q(z),&w\text{ odd},\ p\text{ odd},\\
 U_p(z)U_q(z),&w\text{ odd},\ p\text{ even}.
 \end{cases}
\end{align*}
Then
\begin{align}
 A_p(z)={}&P_p(z)
 +\sum_{\substack{3\le j\le p\\j\text{ odd}}}
       \binom{w-j-1}{q-1}\zeta(j)W_{w-j}(z)
 -\sum_{\substack{3\le j\le q\\j\text{ odd}}}
       \binom{w-j-1}{p-1}\zeta(j)W_{w-j}(z)\notag\\
 &-\frac12\left[\binom{w-1}{q}-\binom{w-1}{p}\right]W_w(z)
 -\boldsymbol{1}_{w\text{ odd}}\frac{(-1)^p}{2}\zeta(w).
 \label{gaussian:eq:parity-A}
\end{align}
Each individual double is recovered by the explicit inverse transform
\begin{equation}\label{gaussian:eq:parity-inverse}
 E_{w-b,b}(z)=\sum_{p=1}^{b}(-1)^{b-p}
                  \binom{w-p-1}{b-p}A_p(z),\qquad 1\le b<w.
\end{equation}
Consequently every imaginary part of even total weight and every real part
of odd total weight in the same-argument double family is an explicit
polynomial in single polylogarithms and ordinary zeta values.
\end{theorem}

\begin{proof}
We supply the analytic identity responsible for the reduction, so that no
rank computation or numerical recognition enters the proof. Define
\[
 T_{p,q}(k)=\sum_{m=1}^{\infty}\frac1{(m+k)^p m^q},\qquad k\ge1.
\]
Partial fractions at $m=0$ and $m=-k$ give
\begin{align}
 \frac1{(m+k)^p m^q}
 ={}&\sum_{j=1}^{q}(-1)^{q-j}
       \binom{w-j-1}{p-1}\frac1{k^{w-j}m^j}\notag\\
 &+(-1)^q\sum_{j=1}^{p}
       \binom{w-j-1}{q-1}\frac1{k^{w-j}(m+k)^j}.
 \label{gaussian:eq:parity-pf}
\end{align}
The two $j=1$ divergences cancel. Summing first to a finite cutoff and
then letting it tend to infinity yields
\begin{align}
 T_{p,q}(k)={}&
 \sum_{j=2}^{q}(-1)^{q-j}\binom{w-j-1}{p-1}
                \frac{\zeta(j)}{k^{w-j}}
 +(-1)^q\sum_{j=2}^{p}\binom{w-j-1}{q-1}
                \frac{\zeta(j)}{k^{w-j}}\notag\\
 &+(-1)^{q+1}\sum_{j=1}^{p}\binom{w-j-1}{q-1}
                \frac{H_k^{(j)}}{k^{w-j}}.
 \label{gaussian:eq:parity-T}
\end{align}
Write
\[
 \mathcal A_p(z)=\sum_{j=1}^{p}\binom{w-j-1}{q-1}F_{w-j,j}(z),
 \qquad C_p=\binom{w-1}{q},
\]
and
\[
 \mathcal U_{p,q}(z)=
 \sum_{j=2}^{q}(-1)^{q-j}\binom{w-j-1}{p-1}
                 \zeta(j)\Li_{w-j}(z)
 +(-1)^q\sum_{j=2}^{p}\binom{w-j-1}{q-1}
                 \zeta(j)\Li_{w-j}(z).
\]
Using $H_k^{(j)}=H_{k-1}^{(j)}+k^{-j}$ and the hockey-stick identity,
equation~\eqref{gaussian:eq:parity-T} becomes
\begin{equation}\label{gaussian:eq:parity-R}
 R_{p,q}(z):=\sum_{k\ge1}z^kT_{p,q}(k)
 =\mathcal U_{p,q}(z)+(-1)^{q+1}
        \bigl(\mathcal A_p(z)+C_p\Li_w(z)\bigr).
\end{equation}

There are two product identities. The familiar same-argument shuffle is
\begin{equation}\label{gaussian:eq:parity-shuffle}
 \Li_p(z)\Li_q(z)=\mathcal A_p(z)+\mathcal A_q(z).
\end{equation}
The Fourier-convolution identity is
\begin{equation}\label{gaussian:eq:parity-convolution}
 \Li_p(z)\Li_q(z^{-1})
       =\zeta(w)+R_{p,q}(z)+R_{q,p}(z^{-1}).
\end{equation}
For completeness, this remains valid when $p=1$ or $q=1$. Each
$\Li_p(e^{i\theta})$ belongs to $L^2$ because $\sum n^{-2p}<\infty$.
The Fourier coefficients of their product are therefore the absolutely
convergent convolutions, equal to $\zeta(w)$ at index zero and to
$T_{p,q}(k)$ and $T_{q,p}(k)$ at positive and negative indices. Both
sequences $T_{p,q}(k)$ decrease to zero. Dirichlet's test gives convergence
of the two Fourier series, uniform on closed subarcs away from $z=1$.
Abel--Poisson summation of the $L^1$ product converges to the continuous
product at every such point, proving~\eqref{gaussian:eq:parity-convolution} there.
Thus the derivation has never multiplied divergent zeta series.

If $w$ is even, take imaginary parts in
\eqref{gaussian:eq:parity-shuffle}--\eqref{gaussian:eq:parity-convolution}. Since $p$ and $q$
have the same parity, these give the sum and difference of $A_p,A_q$.
Moreover,
\begin{align*}
 &\operatorname{Im}\bigl(\mathcal U_{p,q}(z)+
                         \mathcal U_{q,p}(z^{-1})\bigr)\\
 &\quad=2(-1)^p\left[
  \sum_{\substack{3\le j\le p\\j\text{ odd}}}
       \binom{w-j-1}{q-1}\zeta(j)V_{w-j}(z)
 -\sum_{\substack{3\le j\le q\\j\text{ odd}}}
       \binom{w-j-1}{p-1}\zeta(j)V_{w-j}(z)\right].
\end{align*}
Solving the two elementary equations gives~\eqref{gaussian:eq:parity-A} for even
$w$. If $w$ is odd, take real parts. Now $p,q$ have opposite parity and
the same two equations again give the sum and difference. The analogous
last display has $-2(-1)^p$ in place of $2(-1)^p$, with $U$ in place of
$V$. The constant term $\zeta(w)$ contributes
$-(-1)^p\zeta(w)/2$. This proves~\eqref{gaussian:eq:parity-A} in the odd case too.

Finally the coefficient matrix in the definition of $A_p$ is lower
triangular with diagonal one. Its inverse follows from
\[
 \binom{w-j-1}{p-j}\binom{w-p-1}{b-p}
 =\binom{w-j-1}{b-j}\binom{b-j}{p-j}
\]
and $\sum_{p=j}^{b}(-1)^{b-p}\binom{b-j}{p-j}=\delta_{b,j}$.
This proves~\eqref{gaussian:eq:parity-inverse}.
\end{proof}

\subsection{Consequences for Gaussian and Eisenstein values}
\label{gaussian:subsec:gaussian-doubles}

At $z=i$ we have $V_n(i)=\beta(n)$,
$U_1(i)=-\log2/2$, and
\[
 U_n(i)=-2^{-n}(1-2^{1-n})\zeta(n),\qquad n\ge2.
\]
The odd beta values are rational multiples of powers of $\pi$.
Thus Theorem~\ref{gaussian:thm:binomial-parity} supplies an explicit all-weight
reduction to the Gaussian single-value algebra in the appropriate parity.
It proves all five previously numerical weight-six equations:
\begin{align}
\label{gauss:eq:g51}
2048\operatorname{Im}\Li_{5,1}(i,1)
 &=-64\pi^3\zeta(3)-527\pi\zeta(5)+4096\beta(6),\\
\label{gauss:eq:g42}
1536\operatorname{Im}\Li_{4,2}(i,1)
 &=96\pi^3\zeta(3)-32\pi^2\beta(4)+1581\pi\zeta(5)-8448\beta(6),\\
\label{gauss:eq:g33}
1024\operatorname{Im}\Li_{3,3}(i,1)
 &=-3\pi^3\zeta(3)+64\pi^2\beta(4)-1581\pi\zeta(5)+4608\beta(6),\\
\label{gauss:eq:g24}
23040\operatorname{Im}\Li_{2,4}(i,1)
 &=-14\pi^4\beta(2)+135\pi^3\zeta(3)-1440\pi^2\beta(4)\\
 &\hspace{17mm}+23715\pi\zeta(5)-69120\beta(6),\\
\label{gauss:eq:g15}
92160\operatorname{Im}\Li_{1,5}(i,1)
 &=-150\pi^5\log2+56\pi^4\beta(2)-270\pi^3\zeta(3)\\
 &\hspace{17mm}+1920\pi^2\beta(4)-675\pi\zeta(5).
\end{align}
An especially simple all-order subfamily is
\begin{equation}\label{gaussian:eq:leading-index-family}
 \operatorname{Im}\Li_{2m-1,1}(z,1)
  =(m-1)V_{2m}(z)+U_{2m-1}(z)V_1(z)
   -\sum_{k=1}^{m-1}\zeta(2k+1)V_{2m-2k-1}(z),
 \qquad m\ge1.
\end{equation}
For example, this gives the further proved weight-eight identity
\[
 \operatorname{Im}\Li_{7,1}(i,1)
 =3\beta(8)-\frac{5\pi^5}{1536}\zeta(3)
             -\frac{\pi^3}{32}\zeta(5)
             -\frac{8255\pi}{32768}\zeta(7).
\]
At $z=\rho=e^{2\pi i/3}$ write $C_n=\operatorname{Im}\Li_n(\rho)$.
The same theorem applies unchanged; in particular
\[
 \operatorname{Im}\Li_{5,1}(\rho,1)
 =2C_6-\frac{2\pi^3}{81}\zeta(3)-\frac{121\pi}{486}\zeta(5),
\qquad
 \operatorname{Im}\Li_{3,1}(\rho,1)
 =C_4-\frac{13\pi}{54}\zeta(3).
\]
These statements assert reduction identities, not numerical or motivic
independence of the remaining single values.

\section{Classical reductions for every index with one two}
\label{gaussian:sec:one-two}

This section uses the one-variable abbreviation
$\Li_{k_1,\ldots,k_d}(z)=\Li_{k_1,\ldots,k_d}(z,1,\ldots,1)$.
Let $\{1\}^{r}$ denote $r$ consecutive ones. Work first on $0<z<1$, and
then continue analytically along paths avoiding the cuts of the displayed
functions. In particular, the formulas hold at $z=i$ by continuation
through the upper half-plane, with the principal values used below.

\begin{theorem}[One two at arbitrary depth]\label{gaussian:thm:one-two}
For $d\ge1$, putting $a=\log(1-z)$ gives
\begin{equation}\label{gaussian:eq:height-one-general}
 \Li_{2,\{1\}^{d-1}}(z)
 =\zeta(d+1)+\frac{(-1)^d a^d\log z}{d!}
   +\sum_{k=0}^{d-1}\frac{(-1)^{k+1}a^k}{k!}
                       \Li_{d+1-k}(1-z).
\end{equation}
For all $r,s\ge0$, put $\alpha=-\log(1-z)$. Then
\begin{equation}\label{gaussian:eq:one-two-any-position}
 \Li_{\{1\}^{r},2,\{1\}^{s}}(z)
 =\sum_{j=0}^{r}(-1)^j\binom{s+j+1}{j}
      \frac{\alpha^{r-j}}{(r-j)!}
                     \Li_{2,\{1\}^{s+j}}(z).
\end{equation}
Consequently every value with exactly one index equal to two and all other
indices equal to one reduces to classical polylogarithms and logarithms.
\end{theorem}

The height-one functional equation is classical; a modern treatment and
arithmetic analogue appear in \cite{gaussian:NakamuraShiraishi}. The proof below
also records the shuffle reduction when the two occupies any position.

\begin{proof}
The iterated integral for $d$ ones is
$\Li_{\{1\}^{d}}(z)=\alpha^d/d!$. Hence
\[
 \Li_{2,\{1\}^{d-1}}(z)
 =\frac{(-1)^d}{d!}\int_0^z\frac{\log^d(1-t)}{t}\,dt.
\]
Set $s=1-t$ and integrate by parts $d$ times. The antiderivative
\[
 -\log^d s\log(1-s)
  +\sum_{j=1}^{d}(-1)^j\frac{d!}{(d-j)!}
                       \log^{d-j}s\Li_{j+1}(s)
\]
has derivative $\log^d(s)/(1-s)$. At $s=1$ its limiting value is
$(-1)^d d!\zeta(d+1)$. This gives~\eqref{gaussian:eq:height-one-general} including
the constant term and its sign.

For the second identity, locate the unique zero letter in the iterated
integral word $1^r0\,1^{s+1}$. Its variable is $t$. Integrating the
$r$ outer and $s+1$ inner identical letters gives
\[
 \Li_{\{1\}^{r},2,\{1\}^{s}}(z)
 =\frac1{r!(s+1)!}\int_0^z
    \bigl(\alpha(z)-\alpha(t)\bigr)^r\alpha(t)^{s+1}\frac{dt}{t}.
\]
Expanding the first factor and using the preceding integral formula
proves~\eqref{gaussian:eq:one-two-any-position}. Analytic continuation then gives
the asserted complex values.
\end{proof}

\subsection{Proof of all three weight-four Gaussian triple formulas}
\label{gaussian:subsec:gaussian-triples}

Put $L=\log2$, $G=\beta(2)$, $r=(1+i)/2$, and
$\lambda_n=\operatorname{Im}\Li_n(r)$,
$\mu_n=\operatorname{Re}\Li_n(r)$. The elementary dilogarithm reflection,
Landen transformation, and trilogarithm reflection give
\begin{align}
 \mu_2&=\frac{5\pi^2}{96}-\frac{L^2}{8},&
 \lambda_2&=G-\frac{\pi L}{8},\notag\\
 \mu_3&=\frac{35}{64}\zeta(3)-\frac{5\pi^2L}{192}
                       +\frac{L^3}{48}.
 \label{gaussian:eq:halfpoint-auxiliary}
\end{align}
For clarity, the trilogarithm identity used for the last formula is
\begin{align*}
 &\Li_3(z)+\Li_3(1-z)+\Li_3\!\left(-\frac{z}{1-z}\right)\\
 &\quad=\zeta(3)+\zeta(2)\log(1-z)
  -\frac12\log z\log^2(1-z)+\frac16\log^3(1-z).
\end{align*}
It follows by differentiation using the dilogarithm identities and fixing
the constant as $z\to0$; it is then continued to $z=r$, where
$1-r=\bar r$ and $-r/(1-r)=-i$. Taking real parts proves the displayed
$\mu_3$ evaluation. The first two formulas follow respectively from
$\Li_2(r)+\Li_2(\bar r)=\zeta(2)-\log r\log\bar r$ and
$\Li_2(i)+\Li_2(\bar r)=-\tfrac12\log^2(1-i)$.

To evaluate the height-one formulas at $i$, write
\[
 a=\log(1-i)=\frac L2-\frac{\pi i}{4},\qquad
 c=\log(-(1-i))=\frac L2+\frac{3\pi i}{4}.
\]
Since $(1-i)^{-1}=r$, the needed principal-branch inversion formulas are
\begin{align*}
 \Li_2(1-i)&=-\Li_2(r)-\frac{\pi^2}{6}-\frac{c^2}{2},\\
 \Li_3(1-i)&=\Li_3(r)-\frac{\pi^2c}{6}-\frac{c^3}{6},\\
 \Li_4(1-i)&=-\Li_4(r)-\frac{7\pi^4}{360}
                         -\frac{\pi^2c^2}{12}-\frac{c^4}{24}.
\end{align*}
Substitution into~\eqref{gaussian:eq:height-one-general} with $d=2,3$, using
$\log i=\pi i/2$, yields
\begin{align*}
 \operatorname{Re}\Li_{2,1}(i)&=\frac{29}{64}\zeta(3)-\frac{\pi G}{4},\\
 \operatorname{Im}\Li_{2,1}(i)&=-\frac{GL}{2}
                   +\frac{\pi L^2}{32}-\lambda_3+\frac{3\pi^3}{128},\\
 \operatorname{Im}\Li_{2,1,1}(i)&=
 \lambda_4+\frac{L\lambda_3}{2}
 -\frac{G(\pi^2-4L^2)}{32}
 -\frac{\pi(8L^3+105\zeta(3))}{768}.
\end{align*}
The real-part evaluation of $\Li_{2,1}(i)$ is also printed in
Panzer \cite{Panzer2017}, after reversing his index convention.
The remaining two identities follow from exact shuffle relations:
\begin{align*}
 \Li_1(z)\Li_{2,1}(z)
    &=\Li_{1,2,1}(z)+3\Li_{2,1,1}(z),\\
 \Li_2(z)\Li_{1,1}(z)
    &=3\Li_{2,1,1}(z)+2\Li_{1,2,1}(z)+\Li_{1,1,2}(z).
\end{align*}
Using $\Li_1(i)=-a$, $\Li_{1,1}(i)=a^2/2$, and
$\Li_2(i)=-\pi^2/48+iG$, one obtains
\begin{align*}
 \operatorname{Im}\Li_{1,2,1}(i)
 &=-3\lambda_4-L\lambda_3
   +\frac{G(\pi^2-4L^2)}{32}
   -\frac{3\pi^3L}{256}
   +\frac{\pi(2L^3+67\zeta(3))}{128},\\
 \operatorname{Im}\Li_{1,1,2}(i)
 &=3\lambda_4+\frac{L\lambda_3}{2}
   +\frac{5\pi^3L}{192}-\frac{163\pi\zeta(3)}{256}.
\end{align*}
These are exactly the three numerical formulas in the original experimental draft,
now proved. The arbitrary-depth theorem proves additional such identities
without a new integer-relation search. It does not address the independent
status of the weight-six triples $(3,1,2),(3,2,1),(4,1,1)$, which contain
more than one zero letter and lie outside this reduction mechanism.

\section{Classical reductions with one index equal to two}
\label{recip:one2:sec:general}

This section settles the three Gaussian weight-four triple identities
printed in Chapter~4 of the source manuscript.  The proof gives a uniform
formula for every multiple polylogarithm whose indices consist of one
$2$ and any number of $1$'s.  The general phenomenon is classical:
Borwein and Straub give the complementary-argument formula for
$\Li_{2,\{1\}^{p-1}}$ in \cite[Eq.~(33)]{gaussian:BorweinStraub2015}, including
the weight-four instance in their Eq.~(34).  The results below reorganize
that reduction at the reciprocal complementary argument, make the
dependence on the position of the $2$ explicit, and prove the specific
Gaussian formulas in the manuscript.  We do not claim that classical
reducibility of this family is a new phenomenon.

Throughout this section, a single displayed argument abbreviates the
remaining arguments $1$:
\[
 \Li_{s_1,\ldots,s_d}(z)
 :=\Li_{s_1,\ldots,s_d}(z,1,\ldots,1).
\]
The defining series has strictly decreasing summation indices.  Set
\[
 U_{a,b}(z)=\Li_{\{1\}^{a},2,\{1\}^{b}}(z),\qquad
 H_p(z)=\Li_{2,\{1\}^{p-1}}(z),
\]
where $a,b\ge0$ and $p\ge1$.  The weight of $U_{a,b}$ is
$N=a+b+2$, and its series depth is $N-1$.

\subsection{A logarithmic moment and its branch conventions}

We use principal branches and first work on
\[
 \mathcal D=\mathbb C\setminus[0,\infty),\qquad
 w=\frac1{1-z},\qquad L=\Log w=-\Log(1-z).
\]
On this domain, $w$ avoids the standard polylogarithm cut $[1,\infty)$
and the logarithm cut $(-\infty,0]$.  All terms in the formulas below
are therefore single-valued analytic functions.  The formulas also
extend to the usual slit plane by consistent continuation; no such
extension is needed at $z=i$ or at a nonreal root of unity.

\begin{lemma}[Logarithmic moment]
\label{recip:one2:lem:moment}
Let $q(t)=-\Log(1-t)$.  For an integration path from $0$ to $z$ avoiding
$\{0,1\}$ apart from its initial point $0$, with the branches
continued along that path,
\begin{equation}
 U_{a,b}(z)=\frac1{a!(b+1)!}
 \int_0^z\frac{(q(z)-q(t))^a q(t)^{b+1}}{t}\,dt.
 \label{recip:one2:eq:moment}
\end{equation}
Consequently,
\begin{equation}
 U_{a,b}(z)=\sum_{j=0}^a(-1)^j
 \binom{b+j+1}{j}\frac{q(z)^{a-j}}{(a-j)!}\,H_{b+j+1}(z).
 \label{recip:one2:eq:triangular}
\end{equation}
\end{lemma}
\begin{proof}
The iterated-integral word of $U_{a,b}$ is
$1^a0\,1^{b+1}$, with forms $dt/(1-t)$ and $dt/t$ attached to $1$ and
$0$, respectively.  The innermost $b+1$ identical forms contribute
$q(t)^{b+1}/(b+1)!$.  The $a$ identical forms between $t$ and $z$
contribute $(q(z)-q(t))^a/a!$.  Integrating the single remaining $0$
form proves \eqref{recip:one2:eq:moment}.  Equivalently, one may verify the
same formula by differentiating successively with respect to $z$ and
using the zero initial values at $z=0$.  Expanding the binomial and
using
\[
 H_p(z)=\frac1{p!}\int_0^z\frac{q(t)^p}{t}\,dt
\]
proves \eqref{recip:one2:eq:triangular}.
\end{proof}

\begin{theorem}[Explicit reduction for every position of the $2$]
\label{recip:one2:thm:closed}
For $a,b\ge0$, $N=a+b+2$, and $z\in\mathcal D$,
\begin{align}
 U_{a,b}(z)
 &=\sum_{k=a+1}^{N}(-1)^{a+k}\binom{k-1}{a}
       \frac{L^{N-k}}{(N-k)!}\Li_k(w)
       -\frac{L^N}{N!}\notag\\
 &\quad+(-1)^{b+1}\sum_{j=0}^{a}\binom{b+j+1}{j}
       \frac{L^{a-j}}{(a-j)!}\zeta(b+j+2).
 \label{recip:one2:eq:closed}
\end{align}
In particular,
\begin{equation}
 H_p(z)=(-1)^p\zeta(p+1)-\frac{L^{p+1}}{(p+1)!}
 +\sum_{k=1}^{p+1}(-1)^k
       \frac{L^{p+1-k}}{(p+1-k)!}\Li_k(w).
 \label{recip:one2:eq:height}
\end{equation}
\end{theorem}
\begin{proof}
We first prove \eqref{recip:one2:eq:height} for real $z<0$, so that
$0<w<1$ and every intermediate integral is real.  Substitute
$v=1/(1-t)$ in the logarithmic moment.  Then
\[
 \frac{dt}{t}=\frac{dv}{v(v-1)},\qquad q(t)=\log v,
\]
and
\[
 H_p(z)=\frac1{p!}\int_1^w\frac{\log^p v}{v(v-1)}\,dv.
\]
The differential identity
\begin{equation}
 \frac{d}{dv}\left[
 \sum_{k=1}^{p+1}\frac{(-1)^{k-1}p!}{(p+1-k)!}
       \log^{p+1-k}v\,\Li_k(v)\right]
       =\frac{\log^p v}{1-v}
 \label{recip:one2:eq:primitive}
\end{equation}
follows by cancellation of adjacent terms, using
$\Li_1(v)=-\log(1-v)$ and $\Li'_k(v)=\Li_{k-1}(v)/v$ for $k\ge2$.
Since
$1/[v(v-1)]=-1/v-1/(1-v)$, integration gives
\eqref{recip:one2:eq:height}.  At $v=1$, every term in
\eqref{recip:one2:eq:primitive} with a positive power of $\log v$ tends to
zero.  This includes the apparently singular term
$\log^p v\,\Li_1(v)$, which is
$O((v-1)^p\log(1-v))$ and tends to zero because $p\ge1$.
The sole nonzero endpoint term is $(-1)^pp!\zeta(p+1)$.
Thus the constant in \eqref{recip:one2:eq:height} is fixed, not inferred
from a differential identity alone.  Analytic continuation on the
connected domain $\mathcal D$ proves that identity throughout the
stated domain.

Insert \eqref{recip:one2:eq:height} in
\eqref{recip:one2:eq:triangular}, using $q(z)=L$.  The zeta terms immediately
give the second line of \eqref{recip:one2:eq:closed}.  The coefficient of
$(-1)^k L^{N-k}\Li_k(w)$ is
\[
 \sum_{j=0}^a\frac{(-1)^j\binom{b+j+1}{j}}
                    {(a-j)!(b+j+2-k)!}
 =\frac{(-1)^a\binom{k-1}{a}}{(N-k)!}.
\]
Terms containing a factorial of a negative integer are omitted.  To
verify this finite identity, write its left side as
\[
 \frac{(k-1)!}{(b+1)!a!}
 \sum_{j=0}^a(-1)^j\binom aj\binom{b+j+1}{k-1}
\]
and extract the coefficient of $x^{k-1}$ in
\[
 (1+x)^{b+1}\sum_{j=0}^a(-1)^j\binom aj(1+x)^j
   =(-x)^a(1+x)^{b+1}.
\]
The coefficients vanish when $k\le a$, explaining the lower endpoint
$a+1$ in the sum.  Finally, the coefficient of $-L^N$ equals
\begin{align*}
 \sum_{j=0}^a\frac{(-1)^j\binom{b+j+1}{j}}
                    {(a-j)!(b+j+2)!}
 &=\frac1{a!(b+1)!}\int_0^1t^{b+1}(1-t)^a\,dt\\
 &=\frac1{N!}.
\end{align*}
This proves the complete formula.
\end{proof}

For the Gaussian endpoint the path itself can be inspected explicitly:
\[
 z=ix\quad(0\le x\le1),\qquad
 w(x)=\frac{1+ix}{1+x^2}.
\]
For $x>0$ this path lies in the upper half-plane and does not meet the
polylogarithm cut.  Its initial point is $w(0)=1$, where the endpoint
limits in the proof apply, and its final point is $(1+i)/2$.  This
verifies the branch choices directly, without using an inversion
formula with an unspecified logarithmic correction.

The boundary series at $z=i$ agree with these analytic values.
Indeed the finite nested harmonic sums in a coefficient satisfy a
bound $O((1+\log n)^{d-1})$.  After division by $n^{s_1}$ their
coefficients tend to zero and have summable successive differences.
Summation by parts against the bounded partial sums of $i^n$ proves
convergence, even when $s_1=1$; Abel's theorem then identifies the
series with the radial analytic limit.

\begin{corollary}[A short formula when the $2$ is last]
\label{recip:one2:cor:last}
For every $N\ge2$,
\begin{align}
 U_{N-2,0}(z)
 &=(N-1)\Li_N(w)-L\Li_{N-1}(w)-\frac{L^N}{N!}\notag\\
 &\quad-\sum_{m=0}^{N-2}(N-1-m)\zeta(N-m)\frac{L^m}{m!}.
 \label{recip:one2:eq:last}
\end{align}
Thus only two nonconstant classical polylogarithms are needed in this
subfamily, regardless of series depth.
\end{corollary}
\begin{proof}
Set $a=N-2$, $b=0$ in \eqref{recip:one2:eq:closed} and reindex the zeta sum.
\end{proof}

\begin{corollary}[The top two series-depth layers]
\label{recip:one2:cor:toplayers}
Every single-endpoint multiple polylogarithm of weight $N\ge2$ and
series depth at least $N-1$ has an explicit expression in classical
polylogarithms at $(1-z)^{-1}$, zeta values, and logarithms, on the
domain $\mathcal D$.
\end{corollary}
\begin{proof}
Positive indices with sum $N$ and length $N$ are all $1$, and
$\Li_{\{1\}^N}(z)=L^N/N!$. Those with length $N-1$ have exactly one
index $2$ and all other indices $1$, so Theorem~\ref{recip:one2:thm:closed}
applies. These exhaust the possibilities.
\end{proof}

This statement concerns the literal depth of the defining series.
It does not identify that depth with a motivic filtration or prove a
minimal depth for the resulting numerical constants.

\subsection{One possible new direction at each weight}

Put
\[
 \ell=\log2,\qquad w=\frac{1+i}{2},\qquad
 \mu_k=\Re\Li_k(w),\qquad\lambda_k=\Im\Li_k(w).
\]
For $N\ge2$, let $\mathcal P_N$ be the rational span of products of
at least two positive-weight atoms selected from
\[
 \pi,\ell,\quad \zeta(k),\mu_k,\lambda_k\quad(2\le k<N),
\]
whose assigned weights sum to $N$; the assigned weights are $1$ for
$\pi,\ell$ and $k$ for the atoms with subscript or argument $k$.
This definition is a concrete span of numerical products.  It does
not assume that numerical periods form a direct sum by weight.

\begin{corollary}[Signed binomial coefficients modulo products]
\label{recip:one2:cor:primitive}
For $a+b+2=N$, set
$c_{N,a}=(-1)^{N+a}\binom{N-1}{a}$.  Then
\begin{align}
 U_{a,b}(i)&\equiv c_{N,a}\bigl(\Li_N(w)-\zeta(N)\bigr)
        \pmod{\mathcal P_N+i\mathcal P_N},\label{recip:one2:eq:complexprimitive}\\
 \Re U_{a,b}(i)&\equiv c_{N,a}\bigl(\mu_N-\zeta(N)\bigr)
        \pmod{\mathcal P_N},\label{recip:one2:eq:realprimitive}\\
 \Im U_{a,b}(i)&\equiv c_{N,a}\lambda_N
        \pmod{\mathcal P_N}.\label{recip:one2:eq:imagprimitive}
\end{align}
In particular, the imaginary parts of this entire family contribute
at most one direction after products of the stated lower-weight
atoms have been factored out.
\end{corollary}
\begin{proof}
At $z=i$, $L=-\ell/2+i\pi/4$, and
$\Li_1(w)=-\overline L$.  In \eqref{recip:one2:eq:closed}, every term is a
product of lower-weight atoms except the term $k=N$ and the term
$j=a$ in the zeta sum.  Their respective coefficients are
$c_{N,a}$ and $-c_{N,a}$.  Taking real and imaginary parts proves the
statements.
\end{proof}

The coefficients of the highest-weight half-point imaginary value are
\[
 \begin{array}{c|l}
 N&c_{N,0},\ldots,c_{N,N-2}\\\hline
 4&1,-3,3\\
 5&-1,4,-6,4\\
 6&1,-5,10,-10,5.
 \end{array}
\]
The corollary is an upper bound, and proves neither that $\lambda_N$
is nonzero in the quotient nor that it is independent of constants
from other argument families.  It also does not identify series
depth with an additive motivic depth filtration.

\subsection{A sixth-root specialization in every position}
\label{recip:one2:sec:sixthroot}

The special role of sixth roots in relations for Nielsen polylogarithms
is studied by Borwein and Straub~\cite{gaussian:BorweinStraub2015}.  The preceding
position formula gives a direct specialization in that established
setting, including an explicit rational coefficient for one component
at every weight.  Throughout this subsection,
\[
 \omega=e^{i\pi/3},\qquad N=a+b+2,\qquad a,b\ge0,
\]
and all logarithms and polylogarithms use their principal branches.

\begin{corollary}[Sixth-root closure and a pure parity component]
\label{recip:one2:cor:sixthroot}
At $z=\omega$, formula~\eqref{recip:one2:eq:closed} reduces
$U_{a,b}(\omega)$ to classical values $\Li_k(\omega)$, zeta values,
and powers of $i\pi$, since
\begin{equation}
 \frac1{1-\omega}=\omega,\qquad
 -\Log(1-\omega)=\Log\omega=\frac{i\pi}{3}.
 \label{recip:one2:eq:sixthfixed}
\end{equation}
Moreover,
\begin{equation}
 U_{a,b}(\omega)+(-1)^N\overline{U_{a,b}(\omega)}
       =R_{N,a}(i\pi)^N,\qquad R_{N,a}\in\mathbb Q.
 \label{recip:one2:eq:sixthparity}
\end{equation}
Thus $\Re U_{a,b}(\omega)$ is a rational multiple of $\pi^N$
when $N$ is even, and $\Im U_{a,b}(\omega)$ is a rational multiple
of $\pi^N$ when $N$ is odd.  Explicitly, with Bernoulli polynomials
defined by
\[
 \frac{t e^{xt}}{e^t-1}
    =\sum_{k\ge0}B_k(x)\frac{t^k}{k!},\qquad B_k=B_k(0),
\]
one may take
\begin{align}
 R_{N,a}
 &=-\frac{2}{3^N N!}
 -\sum_{k=a+1}^{N}
   \frac{(-1)^{a+k}\binom{k-1}{a}\,2^k B_k(1/6)}
        {3^{N-k}(N-k)!k!}\notag\\
 &\quad+(-1)^{N-a}
   \sum_{\substack{N-a\le m\le N\\m\ \mathrm{even}}}
   \frac{\binom{m-1}{N-a-1}\,2^m B_m}
        {3^{N-m}(N-m)!m!}.
 \label{recip:one2:eq:sixthcoefficient}
\end{align}
The same formulas at $\overline\omega$ follow by complex conjugation.
\end{corollary}
\begin{proof}
The identities in \eqref{recip:one2:eq:sixthfixed} follow from
$1-\omega=\overline\omega$.  Along the radial segment
$z=r\omega$, $0\le r\le1$, the real part of $1-z$ is at least
$1/2$, so the indicated logarithm has no branch crossing.
The transformed path lies in the upper half-plane for $r>0$,
starts at $1$, and ends at $\omega$; the endpoint limits already
used to prove \eqref{recip:one2:eq:closed} apply.  Consequently that
formula has exactly the stated principal values.  The ordinary
boundary multiple series agree with their radial limits by the
same summation-by-parts argument used at $z=i$.

Put $T=i\pi$ and $L=T/3$.  The classical Bernoulli Fourier identity is
\begin{equation}
 \Li_k(e^{2\pi i x})+(-1)^k\Li_k(e^{-2\pi i x})
       =-\frac{(2\pi i)^k}{k!}B_k(x),
 \qquad 0<x<1,\quad k\ge1.
 \label{recip:one2:eq:bernoullifourier}
\end{equation}
For $k=1$, this follows directly from $\Li_1(z)=-\Log(1-z)$
on the stated arc.  Successive integration, with the zero average
of both sides fixing the constants, proves every higher case.
Since $\overline L=-L$, the parity projection of a classical term
in \eqref{recip:one2:eq:closed} is
\begin{align*}
 &L^{N-k}\Li_k(\omega)
       +(-1)^N\overline{L^{N-k}\Li_k(\omega)}\\
 &\hspace{1cm}=L^{N-k}
    \bigl(\Li_k(\omega)+(-1)^k\Li_k(\overline\omega)\bigr)
   =-\frac{2^k B_k(1/6)}{3^{N-k}k!}\,T^N.
\end{align*}
The polynomial term $-L^N/N!$ contributes
$-2T^N/(3^N N!)$.  A zeta term with zeta weight $m$ contributes
zero when $m$ is odd.  For even $m\ge2$, use
\[
 \frac{2\zeta(m)}{(i\pi)^m}=-\frac{2^m B_m}{m!}.
\]
In the last sum of \eqref{recip:one2:eq:closed}, write
$m=b+j+2=N-a+j$.  Its binomial coefficient becomes
$\binom{m-1}{N-a-1}$, while its sign after the even-zeta evaluation
is $(-1)^b=(-1)^{N-a}$.  Combining these three contributions proves
\eqref{recip:one2:eq:sixthcoefficient} and hence
\eqref{recip:one2:eq:sixthparity}.
\end{proof}

There is a precise geometric reason that this specialization closes
on a unit-circle argument.  If $|z|=1$ and $z\ne1$, then
\[
 \left|\frac1{1-z}\right|=1
 \quad\Longleftrightarrow\quad
 |1-z|^2=2-2\Re z=1
 \quad\Longleftrightarrow\quad z\in\{\omega,\overline\omega\}.
\]
Both of these points are fixed by $z\mapsto(1-z)^{-1}$.
This classifies the unit-circle endpoint geometry of this particular
transformation.

\begin{corollary}[The two weight-three values]
\label{recip:one2:ex:sixththree}
Define the ordinary Clausen value by its absolutely convergent series
\[
 C=\operatorname{Cl}_2(\pi/3)
   =\sum_{n\ge1}\frac{\sin(n\pi/3)}{n^2}.
\]
Then
\begin{align}
 U_{0,1}(\omega)=\Li_{2,1}(\omega)
  &=\frac23\zeta(3)-\frac\pi3 C+\frac{i\pi^3}{324},
    \label{recip:one2:eq:sixth21}\\
 U_{1,0}(\omega)=\Li_{1,2}(\omega)
  &=-\frac43\zeta(3)+\frac\pi3 C+\frac{i\pi^3}{324}.
    \label{recip:one2:eq:sixth12}
\end{align}
\end{corollary}
\begin{proof}
The classical values needed are
\[
 \Li_1(\omega)=\frac{i\pi}{3},\qquad
 \Li_2(\omega)=\frac{\pi^2}{36}+iC,\qquad
 \Li_3(\omega)=\frac{\zeta(3)}3+\frac{5i\pi^3}{162}.
\]
The real dilogarithm component and imaginary trilogarithm component
follow from \eqref{recip:one2:eq:bernoullifourier}.  The real
trilogarithm component follows by grouping the defining series
into residue classes, or from the distribution identity
\[
 \Li_s(\omega)+\Li_s(\overline\omega)
   =(1-2^{1-s})(1-3^{1-s})\zeta(s),\qquad s>1.
\]
For example, summing over the three cube roots gives
$(3^{1-s}-1)\zeta(s)$ for the two nontrivial cube roots;
applying the two-term distribution formula to the pair of sixth
roots multiplies that sum by $2^{1-s}-1$.
Now \eqref{recip:one2:eq:height} at $p=2$ gives, with $L=i\pi/3$,
\[
 U_{0,1}(\omega)
   =\zeta(3)-\Li_3(\omega)+L\Li_2(\omega)-\frac23L^3.
\]
Substitution yields \eqref{recip:one2:eq:sixth21}.  The shuffle identity
\[
 \Li_1(\omega)\Li_2(\omega)
     =U_{1,0}(\omega)+2U_{0,1}(\omega)
\]
then proves \eqref{recip:one2:eq:sixth12}.  In both positions,
\eqref{recip:one2:eq:sixthcoefficient} gives $R_{3,a}=-1/162$.
\end{proof}

The independent script \texttt{code/independent/one\_two/sixth\_root.py}
computes \eqref{recip:one2:eq:sixthcoefficient} with exact rational
arithmetic and compares the specialized closed formula, its pure
parity component, and the displayed weight-three examples with the
logarithmic moment.  It covers every position of the $2$ at weights
$2$ through $8$, at both $\omega$ and $\overline\omega$, using
$100$ working decimal digits.  Across these $56$ cases, the largest
closed-form versus moment residual is below $6.59\cdot10^{-100}$;
the largest parity-projection residual is below $7.15\cdot10^{-102}$.
The four checks of the two explicit examples and their conjugates
have residual below $2.16\cdot10^{-101}$.  These are numerical checks
of the proved identities, rather than rigorous quadrature error bounds.


\subsection{Further exact identities at weights five and six}

The same calculation gives formulas beyond the three displayed
manuscript rows.  They involve the unavoidable output of this
particular reduction, namely higher classical values at $w$; no
independence claim about those values is implied.

\begin{corollary}[Two higher-weight Gaussian reductions]
\label{recip:one2:cor:56}
With $H_4=\Li_{2,1,1,1}$ and $H_5=\Li_{2,1,1,1,1}$,
\begin{align}
 \Im H_4(i)
 &=-\lambda_5-\frac{\ell\lambda_4}{2}+\frac{\pi\mu_4}{4}
       +\frac{\pi^2-4\ell^2}{32}\lambda_3\notag\\
 &\quad+\frac{G\ell(3\pi^2-4\ell^2)}{192}
       +\frac{35\pi\ell\zeta(3)}{512}
       +\frac{\pi\ell^2(\ell^2-\pi^2)}{384}
       -\frac{17\pi^5}{92160},
       \label{recip:one2:eq:H4}\\
 \Im H_5(i)
 &=\lambda_6+\frac{\ell\lambda_5}{2}-\frac{\pi\mu_5}{4}
       +\frac{4\ell^2-\pi^2}{32}\lambda_4
       -\frac{\pi\ell\mu_4}{8}\notag\\
 &\quad+\frac{\ell(4\ell^2-3\pi^2)}{192}\lambda_3
       +\frac{G(16\ell^4-24\ell^2\pi^2+\pi^4)}{6144}\notag\\
 &\quad+\frac{35\pi(\pi^2-12\ell^2)\zeta(3)}{24576}
       +\frac{\pi\ell^3(5\pi^2-3\ell^2)}{5760}.
       \label{recip:one2:eq:H5}
\end{align}
\end{corollary}
\begin{proof}
Use $p=4$ and $p=5$ in \eqref{recip:one2:eq:height}, substitute
$L=-\ell/2+i\pi/4$ and Lemma~\eqref{gaussian:eq:halfpoint-auxiliary}, and take
imaginary parts.  The accompanying coefficient generator performs
these polynomial operations over $\mathbb Q$ and also emits every
position of the $2$ at weights four, five, and six.
\end{proof}

The real parts can be treated by the same formula.  For example,
\begin{equation}
 \Re\Li_{2,1,1}(i)=\mu_4+\frac{\pi\lambda_3}{4}
 +\frac{\pi G\ell}{8}+\frac{35\ell\zeta(3)}{128}
 +\frac{\ell^4}{384}-\frac{11\pi^2\ell^2}{768}
 -\frac{1249\pi^4}{92160}.
 \label{recip:one2:eq:real211}
\end{equation}
The machine-readable tables include both real and imaginary parts,
so this companion family can be reproduced without manual
transcription.


\section{From one-two reductions to general complementary depth}

For positive indices we abbreviate
\begin{equation}\label{complement:eq:mpl}
\Li_{s_1,\ldots,s_d}(z)
 :=\Li_{s_1,\ldots,s_d}(z,1,\ldots,1)
 =\sum_{n_1>\cdots>n_d\ge1}
 \frac{z^{n_1}}{n_1^{s_1}\cdots n_d^{s_d}},\qquad |z|<1.
\end{equation}
The indices decrease in the summation order, as in the manuscript.
Weight is $w=s_1+\cdots+s_d$ and series depth is $d$.
An empty index denotes the multiplicative identity.
We use convergent multiple zeta values
$\zeta(\boldsymbol s)=\Li_{\boldsymbol s}(1)$ when $s_1\ge2$.

Let
\[
 \omega_0(t)=\frac{dt}{t},\qquad
 \omega_1(t)=\frac{dt}{1-t}.
\]
Words are written \emph{outer letter first}. Thus
\begin{equation}\label{complement:eq:word}
H_{0^{s_1-1}1\cdots0^{s_d-1}1}(z)=\Li_{s_1,\ldots,s_d}(z).
\end{equation}
In particular $H_{01}=\Li_2$, not its negative or a reversed word.
The recursion is
$H_{av}(z)=\int_0^z\omega_a(t)H_v(t)$ whenever this is an ordinary
convergent integral. Pure-zero words are defined by
$H_{0^r}(z)=\log^r z/r!$, and trailing-zero words are obtained by shuffle
regularization. Every word ending in $1$ is convergent at its initial
endpoint and has the ordinary interpretation in \eqref{complement:eq:word}.

Throughout the functional identities we use
\begin{equation}\label{complement:eq:domain}
 D=\C\setminus[0,\infty),\qquad
 q=\frac1{1-z},\qquad L=\log q=-\log(1-z),\qquad z\in D,
\end{equation}
with principal logarithms and principal polylogarithms. The image of $D$
avoids both $(-\infty,0]$ and $[1,\infty)$ in the $q$-plane, so the
hyperlogarithms with possible trailing zeros have unambiguous branches.
The origin is a limiting point, not an included point of $D$.

We first prove formulas on $z<0$, where $0<q<1$ and all logarithms in
$q$ are real. Analytic continuation on the connected slit domain $D$
then proves the stated complex identities. This method avoids an
unstated choice of a value on the polylogarithm cut. It includes both
$z=i$ and $z=\rho=e^{2\pi i/3}$:
\begin{align}
 z=i:&\quad q=\frac{1+i}{2},&L&=-\frac12\log2+\frac{i\pi}{4},\label{complement:eq:gausspoint}\\
 z=\rho:&\quad q=\frac12+\frac{i\sqrt3}{6},&L&=-\frac12\log3+\frac{i\pi}{6}.\label{complement:eq:eispoint}
\end{align}
Their transformed moduli are $2^{-1/2}$ and $3^{-1/2}$ respectively.

\begin{definition}[Depth with logarithmic coefficients]
A term $L^j\Li_{\boldsymbol a}(q)\zeta(\boldsymbol b)$ has
\emph{reduced depth} $\dep(\boldsymbol a)+\dep(\boldsymbol b)$.
The factor $L^j$ is a coefficient, not an additional nested sum.
\end{definition}
This is an explicit convention for the following reduction. It is not
an identification with every motivic depth filtration, and does not
turn arbitrary products of deep factors into depth-one objects.

\section{An explicit complementary-depth theorem}

\subsection{Endpoint transport with convergent constants}
For a word $v$ ending in $0$, write $J_v(q)$ for the iterated integral
from $1$ to $q$ using the same forms and the same outer-first convention.
These integrals converge at $1$: the innermost form is $\omega_0$, which
is regular there; preceding $\omega_1$ integrations do not destroy
integrability. Set $J_\emp=1$.

\begin{lemma}[Connection formula]\label{complement:lem:connection}
For every nonempty word $v$ ending in $0$ and $0<q<1$,
\begin{equation}\label{complement:eq:connection}
 J_v(q)=\sum_{v=uv'}(-1)^{|v'|}H_u(q)H_{\operatorname{rev}(v')}(1).
\end{equation}
The sum includes empty prefixes and suffixes. Every nonempty endpoint
word $\operatorname{rev}(v')$ starts in $0$, so its value at $1$ is
convergent after the usual trailing-zero normalization at $0$.
\end{lemma}
\begin{proof}
We give the endpoint justification rather than silently substituting into
a divergent path identity. Fix a maximum word length $N$ and work in the
truncated noncommutative algebra on $e_0,e_1$. Let
$F(q)=\sum_{|u|\le N}H_u(q)u$. It satisfies
\[
 F'(q)=\left(\frac{e_0}{q}+\frac{e_1}{1-q}\right)F(q).
\]
Put $\ell=-\log(1-q)$. The function $e^{-\ell e_1}F(q)$ has a finite
coefficientwise limit $Z$ as $q\to1^-$: its differential equation has
coefficient $e^{-\ell e_1}e_0e^{\ell e_1}/q$, a finite polynomial in
$\ell$ at each truncated weight, and this is integrable at $1$.
Consequently
$F(q)=e^{\ell e_1}Z+O((1-q)\log^M(1-q))$
coefficientwise, for a sufficiently large finite $M$.

Shuffle identities say that $F$ is group-like; the same holds for $Z$.
The inverse coefficient in a group-like word series is
$(Z^{-1})_{v'}=(-1)^{|v'|}Z_{\operatorname{rev}(v')}$, by the shuffle
antipode, or equivalently by reversal of an iterated path.
The coefficients of $R(q)=F(q)Z^{-1}$ solve the same differential
recursion and tend to zero at $1$ for every nonempty word ending in $0$.
Induction on word length therefore identifies them with the convergent
integrals $J_v(q)$.

If $v'$ is a nonempty suffix of $v$, its reverse begins in $0$.
The corresponding coefficient of $Z$ is simply
$H_{\operatorname{rev}(v')}(1)$: multiplication by
$e^{-\ell e_1}$ only changes words beginning with $1$.
Taking the coefficient of $v$ in $F(q)Z^{-1}$ proves
\eqref{complement:eq:connection}. Since $N$ was arbitrary, this proves every finite
word identity.
\end{proof}

\subsection{The word formula and the depth bound}
For a word $w$ ending in $1$, let $K$ be its number of zeros. Replace
each original $1$ by $0$, and replace each original $0$ independently
by either $0$ or $1$. Denote the resulting $2^K$ words by $\tau_A(w)$,
where $A$ records the zero positions replaced by $1$.
Every $\tau_A(w)$ ends in $0$.

\begin{theorem}[Complementary depth]\label{complement:thm:complement}
For $z\in D$ and every nonempty word $w$ ending in $1$,
\begin{equation}\label{complement:eq:master}
 H_w(z)=(-1)^K\sum_A\ \sum_{\tau_A(w)=uv}
 (-1)^{|v|}H_u(q)H_{\operatorname{rev}(v)}(1).
\end{equation}
After trailing-zero normalization, this is a finite rational linear
combination of terms
\begin{equation}\label{complement:eq:termshape}
 L^j\Li_{\boldsymbol a}(q)\zeta(\boldsymbol b),\qquad
 j+\wt(\boldsymbol a)+\wt(\boldsymbol b)=|w|,
 \quad \dep(\boldsymbol a)+\dep(\boldsymbol b)\le K.
\end{equation}
Nonempty zeta indices are convergent. Thus a one-variable polylogarithm
of weight $w$ and depth $d$ has the stated reduction with $K=w-d$.
\end{theorem}
\begin{proof}
On $z<0$, make the substitution $u=(1-t)^{-1}$ in every integration.
Its pullbacks are
\begin{equation}\label{complement:eq:pullbacks}
 \frac{dt}{t}=-\frac{du}{u}-\frac{du}{1-u},\qquad
 \frac{dt}{1-t}=\frac{du}{u}.
\end{equation}
The initial endpoint $t=0$ becomes $u=1$. Multilinearity gives
$H_w(z)=(-1)^K\sum_A J_{\tau_A(w)}(q)$.
Lemma~\ref{complement:lem:connection} proves \eqref{complement:eq:master}.

It remains to check the claimed normal form. The identity
$H_vH_0=H_{v\sh0}$ removes a trailing zero. If a word ends in exactly
$r$ zeros, its coefficient in the shuffle of its prefix with $0$ is
$r$; all other words in that shuffle have fewer trailing zeros.
Induction therefore expresses each $H_u(q)$ as a rational combination
of $L^jH_{u'}(q)$ with $u'$ empty or ending in $1$.
The operation preserves the number of $1$ letters.

For an endpoint word beginning in $0$, the same normalization preserves
its initial $0$ in every nonzero term with no logarithmic factor.
Terms containing a positive power of $\log1$ vanish. The surviving
nonempty terms begin in $0$ and end in $1$, hence are convergent MZVs.
In each split of $\tau_A(w)$ the total number of $1$ letters in the
two factors is $|A|\le K$. This proves the depth bound and the weight
identity. Finally continue analytically from $z<0$ to $D$.
\end{proof}

\begin{corollary}[An argument-enlarged depth bound]\label{complement:cor:halfb}
For $w\ge2$, every one-variable MPL of weight $w$ can be represented
using nonlogarithmic factors of depth at most $\lfloor w/2\rfloor$,
allowing both the original argument $z$ and $q=(1-z)^{-1}$, MZVs, and
powers of $L$ as coefficients. More precisely, use the original
expression when $d\le w/2$, and Theorem~\ref{complement:thm:complement} otherwise.
The all-one case is $L^w/w!$.
\end{corollary}

This corollary is deliberately not phrased as a theorem about arbitrary
multivariable MPLs, or as a reduction inside an unchanged root-of-unity
basket. At $z=i$ the half-point $q$ is not a root of unity. The theorem
explains an argument-enlarged upper bound; it does not prove the absence
or presence of new cyclotomic periods in a smaller comparison algebra.


For the next family put $F_r(z)=\Li_{2,\ones{r-1}}(z)$, as in the
one-two formulas above. Their reciprocal-complementary form is
\[
 F_r(z)=\sum_{k=1}^{r+1}\frac{(-1)^kL^{r+1-k}}{(r+1-k)!}\Li_k(q)
 +(-1)^r\zeta(r+1)-\frac{L^{r+1}}{(r+1)!}.
\]
\section{A double-polylogarithm formula for an all-weight family}

The next formula addresses the first case beyond the classical ladder.
It reduces the height-one family with initial index $3$ to singles and
doubles at $q$, irrespective of its original depth.
Let
\begin{equation}\label{complement:eq:Ck}
 C_1(q)=0,\qquad C_k(q)=\sum_{a=1}^{k-1}\Li_{a,k-a}(q)\quad(k\ge2).
\end{equation}
All MPLs on the right have the one-variable meaning \eqref{complement:eq:mpl}.

\begin{theorem}[Two-zero height-one identity]\label{complement:thm:twozero}
For every $w\ge3$ and $z\in D$,
\begin{align}
\Li_{3,\ones{w-3}}(z)
={}&\frac{L^w}{w!}
 +\sum_{k=1}^{w}\frac{(-1)^kL^{w-k}}{(w-k)!}
   \bigl((k-2)\Li_k(q)+C_k(q)\bigr)\notag\\
&+(-1)^{w-1}\Bigl((L+\Li_1(q))\zeta(w-1)
 +(w-2)\zeta(w)-\zeta(w-1,1)\Bigr).
\label{complement:eq:twozero}
\end{align}
The right side has reduced depth at most two.
\end{theorem}

\begin{proof}
Write $D_q=q\,d/dq=d/dL$. Direct differentiation of the nested series,
followed by continuation, gives
\[
 D_qC_k=C_{k-1}+\frac{q}{1-q}\Li_{k-1}\quad(k\ge2).
\]
Put $V_k=(k-2)\Li_k+C_k$, with $V_1=-\Li_1$.
Then
\[
 D_qV_k=V_{k-1}+\frac{\Li_{k-1}}{1-q}\quad(k\ge2),
 \qquad D_qV_1=-\frac{q}{1-q}.
\]
Differentiating the finite sum in \eqref{complement:eq:twozero} telescopes and yields
\[
 \frac{F_{w-2}(z)-(-1)^{w-2}\zeta(w-1)}{q-1}.
\]
The derivative of the explicit $(L+\Li_1(q))\zeta(w-1)$ term supplies
the missing constant term, because
$D_q(L+\Li_1(q))=1/(1-q)$.
The derivative of the entire right side is therefore
$F_{w-2}(z)/(q-1)$.
The left side has the same derivative, since
$d\Li_{3,\ones{w-3}}(z)/dz=F_{w-2}(z)/z$ and
$dz/z=dL/(q-1)$.

For the endpoint constant, recall the convergent depth-two sum identity
\begin{equation}\label{complement:eq:eulersum}
 \sum_{a=2}^{w-1}\zeta(a,w-a)=\zeta(w),\qquad w\ge3.
\end{equation}
It follows, for example, by subtracting the shuffle and stuffle
expansions of the once-regularized product $\zeta(1)\zeta(w-1)$:
its common divergent term $\zeta(1,w-1)$ cancels, leaving precisely
\eqref{complement:eq:eulersum}. Only one logarithmic divergence is involved, so
no higher regularization correction enters. Equivalently one may use
finite cutoffs and subtract the common harmonic-logarithmic part
before taking the limit.

The convergent colored stuffle identity, for $0<q<1$, gives
\[
 \Li_{1,w-1}(q)-\Li_1(q)\zeta(w-1)
 =-\Li_{w-1,1}(1,q)-\Li_w(q).
\]
The right side tends to $-\zeta(w-1,1)-\zeta(w)$ by dominated
convergence. Combining this with \eqref{complement:eq:eulersum} yields
\[
 \lim_{q\to1^-}\bigl(C_w(q)-\Li_1(q)\zeta(w-1)\bigr)
 =-\zeta(w-1,1).
\]
All positive powers of $L$ multiplying at most logarithmically singular
terms vanish. The remaining constants in \eqref{complement:eq:twozero} cancel,
so its right side tends to zero as $z\to0^-$. The left side does too.
Equality of derivatives and the endpoint limit proves the result;
continue to $D$.
\end{proof}

\begin{remark}[The sum identity and the proof dependency]
Equation~\eqref{complement:eq:eulersum} is the classical depth-two Euler sum
formula, not a new conjecture assumed here. The same two-zero formula
can also be obtained directly from Theorem~\ref{complement:thm:complement} by
expanding its two zero positions and normalizing trailing zeros.
The supplied code verifies the agreement of these two independently
specified formulas through weight twelve.
\end{remark}

At weight six, \eqref{complement:eq:twozero} is a complete depth-two representation
of $\Li_{3,1,1,1}(i)$ and of $\Li_{3,1,1,1}(\rho)$.
In particular a depth-four index does not by itself identify a new
depth-four period in the enlarged comparison algebra.
The endpoint double is harmless as a named convergent constant;
standard Euler reductions can simplify it further. For example,
\[
 \zeta(4,1)=2\zeta(5)-\zeta(2)\zeta(3),\qquad
 \zeta(5,1)=\frac34\zeta(6)-\frac12\zeta(3)^2.
\]
These last simplifications are not required for the theorem.


